\bookchapter{제8장: 커뮤니케이션 — 문자, 전화, 소셜}

\chapterintro{커뮤니케이션은 휴대폰이 탄생한 가장 원초적이고 핵심적인 목적입니다. 지금은 단순한 통화와 문자에서 메신저, 소셜 네트워크, 이메일까지 아우르는 복잡한 생태계로 발전했습니다.}

\booksection{전화 앱의 고급 기능}

최신 휴대폰의 전화 앱은 단순한 다이얼 그 이상입니다. 통화 녹음, 스팸 차단, 스마트 번호 식별 같은 강력한 도구가 통합되어 있습니다.

\booksubsection{통화 녹음 (법적 준수에 주의):}

\begin{itemize}
  \item 용도: 중요한 통화 내용을 기록해 업무 메모, 인터뷰, 고객 약속, 분쟁 대비 자료로 활용합니다.
  \item 법률 팁: 중국에서는 통화 녹음을 증거로 쓰려면 상대방 동의를 받거나, 자신의 정당한 권리를 보호하는 목적이어야 합니다.
  \item 일방 녹음 vs 쌍방 동의: 법적으로 "상대방 동의 없이" 녹음한 내용이 무효인 것은 아닙니다. 최고인민법원의 관련 사법 해석에 따르면, 녹음 취득 방식이 타인의 정당한 권리를 심각하게 침해하지 않고, 법의 금지 규정을 위반하지 않으며, 공서양속에 어긋나지 않으면 증거로 제출할 수 있습니다. 다만 일부 성급 법원은 민사 재판에서 "몰래 녹음"의 채택 기준을 더 엄격하게 적용하며, 형사 사건에서는 보통 더 엄격한 적법성 심사를 요구합니다.
  \item 지역별 차이: 베이징, 상하이, 광둥 등지의 법원은 일방 녹음을 민사 증거로 비교적 관대하게 인정합니다. 그러나 영업 비밀이나 개인 프라이버시가 걸린 녹음은 어느 성에서든 특별히 신중해야 하며, 프라이버시 침해로 증거에서 제외될 수 있습니다.
  \item 사용 방법: 안드로이드는 보통 통화 화면에 "녹음" 버튼이 있습니다. iOS는 기본 통화 녹음 기능을 제공하지 않으므로 타사 앱이나 녹음 기능이 있는 블루투스 이어폰을 활용해야 합니다.
\end{itemize}

\begin{tcbwarn}{녹음 사용 금지선}
\begin{itemize}
  \item 녹음을 협박, 유포, 공개적 모욕 목적으로 사용해서는 안 됩니다.
  \item 제3자 프라이버시(예: 우연히 타인의 대화가 녹음된 경우)가 포함된 내용은 함부로 퍼뜨리지 마세요.
  \item 중요한 상황에서는 통화 시작 시 "이 통화는 녹음됩니다"라고 구두로 알리는 것이 좋습니다. 규정 준수에도, 증거 효력에도 유리합니다.
\end{itemize}

\end{tcbwarn}

\booksubsection{스팸 전화 대응 심화:}

\begin{itemize}
  \item 원리: 전화 앱이 네트워크의 방대한 번호 데이터베이스를 조회해 수신 시 "스팸 전화", "보이스피싱", 회사 이름 등을 표시합니다.
  \item 켜는 방법: 안드로이드는 전화 앱 > 설정 > 스팸 차단, iOS는 설정 > 전화 > 발신 차단 및 발신자 식별로 들어갑니다. 브랜드별 상세 경로와 권장 설정은 13장 "유용한 시스템 도구"의 브랜드 비교표를 참고하세요.
  \item 주요 번호 표시 플랫폼: 중국 내 번호 표시는 주로 세 곳의 데이터 서비스사가 담당합니다 — Tencent Mobile Manager, 360 Mobile Guard, 전화방(Tapir). 통신사(China Mobile/Unicom/Telecom)의 "가입자 이름" 표시와 "스팸 차단" 기능도 이 데이터를 자주 활용합니다. 휴대폰 제조사(화웨이, 샤오미, OPPO, vivo)의 표시 데이터베이스도 대부분 이 서비스사들과 협력합니다.
  \item 잘못 표시됐다면: 정상 업무 번호가 "배달", "택배", "스팸 전화" 등으로 잘못 표시됐다면 아래 경로로 이의를 신청할 수 있습니다. 전화방(Tapir): 번호 이의신청 플랫폼(haoma.tk.cn)에 접속해 신분증명을 제출하세요.
  \item 360: 360 Mobile Guard 앱의 "번호 이의신청" 기능으로 제출하세요.
  \item Tencent: "Tencent Mobile Manager" 미니프로그램의 "번호 이의신청" 페이지에서 제출하세요.
  \item 3대 통신사: China Mobile 이용자는 "KTFSR" 문자를 10086으로 보내 스팸 차단 기능을 켤 수 있습니다. China Telecom 이용자는 "Tianyi Anti-Spam" 공식 계정을 팔로우하고, China Unicom 이용자는 "Unicom Mobile Manager"에서 설정합니다.
\end{itemize}

스팸 차단 고급: 통신사의 "고빈도 발신 차단"(예: China Mobile의 고빈도 발신 제한)을 켜면 AI 음성 자동 발신도 자동으로 차단됩니다.

휴대폰 기본 차단 주요 기능: 발신자 식별 / 스마트 차단: 번호 데이터베이스를 기준으로 스팸, 사기, 광고 전화를 자동 표시하고 차단합니다.

블랙리스트 / 화이트리스트: 알려진 스팸 번호는 직접 차단하고, 연락처·택배·배달 번호는 화이트리스트에 등록해 오차단을 방지하세요.

문자 키워드 필터: "당첨", "대출", "링크 클릭" 같은 키워드가 포함된 문자는 자동으로 차단함으로 분류됩니다.

iOS 알 수 없는 발신자 음소거: 연락처에 없는 수신 전화는 음성 메일로 조용히 연결되어 방해받지 않습니다. 스팸 전화가 많은 사용자에게 좋지만, 모르는 번호의 중요한 전화를 놓칠 수 있으니 주의하세요.

방해 금지 모드 연동: 야간이나 회의 중에는 연락처 수신만 허용하고 예약 켜기/끄기를 함께 쓰면 수면 질이 바로 나아집니다.

실용 앱 조합: 시스템 기본 차단을 먼저 쓰고, 그 위에 국가 사기 신고 센터와 통신사 스팸 차단을 덧붙이세요. 정리 앱은 Qingzhuo 같은 가벼운 도구를 선택하되, 보안/정리 앱은 한 종류만 설치하면 충분합니다. 자세한 내용은 13장을 참고하세요.

\booksubsection{VoLTE/VoNR가 무엇이고 켜는 방법:}

\begin{itemize}
  \item VoLTE(Voice over LTE): 4G 네트워크 기반 음성 통화 기술로, 통화 중 2G/3G로 떨어지지 않습니다. 통화하면서 인터넷을 쓸 수 있고 음질이 더 선명하며 연결도 빠릅니다.
  \item VoNR(Voice over NR): 5G 기반 음성 통화로, VoLTE의 다음 세대 진화입니다. 지연이 더 낮고 음질이 더 좋으며 고화질 영상 통화도 지원합니다. 5G 휴대폰 + 5G 요금제 + 현지 5G 커버리지가 모두 갖춰져야 쓸 수 있습니다.
  \item 켜는 방법: iOS는 "설정 > 셀룰러 > 셀룰러 데이터 옵션 > 음성 및 데이터"에서 "VoLTE"를 켭니다. 안드로이드는 "설정 > 모바일 네트워크 > 모바일 데이터 > 고급"에서 "VoLTE 고음질 통화"를 찾습니다. 통신사는 보통 기본으로 켜 두는데, 통화 중 네트워크가 끊기거나 인터넷이 안 되면 대부분 VoLTE가 꺼져 있는 경우입니다.
\end{itemize}

\begin{tcbinfo}{알아두면 좋아요}
통화 중 인터넷 속도가 갑자기 느려지거나 끊긴다면 대개 VoLTE/VoNR이 꺼져 2G/3G로 내려간 상태입니다. 켜면 “통화하면서 영상 보기”가 가능해집니다.

\end{tcbinfo}

\booksection{문자 메시지와 iMessage/RCS}

전통적인 문자 메시지는 아직 사라지지 않았습니다. 인증번호, 은행 알림, 공식 정보를 받는 신뢰할 수 있는 채널입니다.

\begin{itemize}
  \item iMessage: 애플 기기 간에 인터넷으로 주고받는 메시지로 파란 말풍선에 표시됩니다. 고화질 사진·영상, 읽음 확인 등을 지원하며 완전 무료이고 종단 간 암호화를 제공합니다.
  \item RCS: 전 세계 통신사와 Google이 밀어붙이는 차세대 커뮤니케이션 프로토콜입니다. iMessage나 WeChat과 비슷한 기능을 전화번호 기반으로 제공합니다.
  \item iMessage와 RCS 상호 운용: iOS 18부터 애플이 RCS 표준을 지원하기 시작했습니다. 이제 iPhone과 안드로이드폰 사이에서 RCS로 고화질 사진, 읽음 확인, 입력 중 표시를 주고받을 수 있어 "초록 말풍선"의 단절감이 줄어듭니다. 단, 양쪽 통신사가 RCS를 지원하고 시스템 버전이 충분히 최신이어야 합니다.
\end{itemize}

\booksubsection{중국 내 RCS 도입 현황:}

\begin{itemize}
  \item China Mobile: 비교적 일찍 RCS를 도입했고, 일부 성에서는 "5G 메시지" 서비스를 지원해 리치 미디어 메시지와 인터랙티브 카드를 보낼 수 있습니다. 단말 지원이 필요하며 대리점에서 개통해야 합니다.
  \item China Telecom: "Telecom 5G 메시지"로 서비스하며, 일부 요금제는 기본 지원됩니다. 기존 이용자는 지정된 문자 명령으로 개통해야 합니다.
  \item China Unicom: 주요 도시에 5G 메시지를 도입했지만 단말 커버리지가 고르지 않아, 많은 안드로이드 기종에서 문자 앱의 "향상된 정보" 또는 "5G 메시지" 스위치를 직접 켜야 합니다.
  \item 실제 체감: 중국 내 RCS 확산은 예상보다 느립니다. WeChat 생태계가 강하고 단말 적합성이 통일되지 않은 탓이 큽니다. 일상 커뮤니케이션은 여전히 WeChat이 중심이고, RCS는 "공식 알림"(은행 알림, 행정 메시지 등) 역할에 더 가깝습니다.
\end{itemize}

\begin{tcbwarn}{인증번호 보안 금지선}
\begin{itemize}
  \item 인증번호를 전달하지 마세요: "인증번호를 읽어 달라", "문자를 전달해 달라"는 사람은 누구든 사기꾼일 수 있습니다. "고객센터", "택배 기사", "공안·검찰·법원"을 자처하는 사람도 포함됩니다. 인증번호 = 당신의 지갑입니다.
  \item 스미싱 문자 주의: "ETC 만료", "의료보험카드 동결", "택배 분실 보상" 등의 문자 속 링크를 경계하세요. 정상 기관은 문자 링크로 은행카드 번호나 비밀번호를 요구하지 않습니다.
  \item 번호 식별: 은행 공식 문자는 보통 95로 시작하는 짧은 번호나 1069+기업 서명에서 옵니다. 일반 11자리 휴대폰 번호에서 온 "은행 알림"은 거의 다 사기입니다.
  \item 2차 인증 백업: 중요한 계정에 묶인 휴대폰 번호는 함부로 바꾸지 마세요. 꼭 바꿔야 한다면 먼저 모든 핵심 계정의 연동을 해제하세요. 새 번호 주인이 문자 인증번호로 당신의 비밀번호를 재설정할 수 있습니다.
\end{itemize}

\end{tcbwarn}

\booksection{WeChat/QQ: 국민 메신저 관리}

WeChat과 QQ는 커뮤니케이션 도구를 넘어 소셜, 결제, 정보 습득을 하나로 묶은 “슈퍼 앱”이 되었습니다.

\booksubsection{채팅 기록 백업 및 이전:}

\begin{enumerate}
  \item PC 또는 Mac에서 WeChat 데스크톱 버전에 로그인합니다.
  \item 왼쪽 아래 메뉴 > [이전 및 백업] > [백업 및 복원]을 클릭합니다.
  \item "채팅 기록을 컴퓨터에 백업"을 선택하고, 휴대폰에서 백업할 채팅을 고릅니다.
  \item 새 휴대폰으로 이전: 기존 휴대폰 WeChat에서 [나] > [설정] > [일반] > [채팅 기록 이전 및 백업] > [이전] > [휴대폰/태블릿으로 이전]을 선택하고 QR을 스캔하면 끝입니다.
\end{enumerate}

\booksubsection{저장 공간 정리 심화 방법:}

\begin{itemize}
  \item 일반 경로: WeChat [나] > [설정] > [일반] > [저장 공간]에서 캐시를 정리하고, 시간·파일 유형별로 채팅 기록을 세밀하게 정리할 수 있습니다.
  \item 채팅 내 정리: 특정 채팅으로 들어가 오른쪽 위 "…" > [채팅 기록] > [저장 공간 관리]에서 그 채팅의 사진·영상·파일만 따로 지우고 텍스트는 남길 수 있습니다.
  \item 원본 사진/영상 정리: 단체 채팅에서 받은 고화질 원본이 공간을 가장 많이 차지합니다. [저장 공간 관리]에서 "원본 사진"을 선택해 일괄 삭제하세요.
  \item 자동 다운로드 끄기: [설정] > [일반] > [사진, 동영상, 파일 및 통화]에서 "자동 다운로드"를 끄면 WiFi에서 단체 파일이 조용히 쌓이는 걸 막을 수 있습니다.
  \item 미니프로그램 캐시: [발견] > [미니프로그램]에서 자주 쓰지 않는 미니프로그램을 길게 눌러 삭제하면 로컬 캐시가 비워집니다.
  \item 채팅 기록 압축 백업: PC WeChat 백업 시 "압축"을 선택하면 수 GB짜리 채팅 기록을 원본의 30\textasciitilde{}50\%로 줄일 수 있습니다.
\end{itemize}

\booksubsection{단체 채팅 관리 팁:}

\begin{itemize}
  \item 공지: 그룹장/관리자가 채팅 오른쪽 위 "…" > [공지]에서 게시하면 전체 멤버에게 강조 알림이 갑니다. 중요한 공지에 적합하며, 새 멤버 입장 시 자동으로 뜹니다.
  \item 그룹 할 일: 그룹장이 고정 메시지를 "할 일"로 지정하면 모든 멤버 머리 위에 빨간 알림이 뜨고, 완료를 누르면 사라집니다. 이벤트 신청, 과제 수집에 좋습니다.
  \item 그룹 이어달리기: 입력창에 '\#이어달리기'를 입력하면 이어달리기를 시작할 수 있어요. 같이 주문하거나 신청자를 모을 때 딱이에요. 이어달리기 결과는 자동으로 단체 채팅 상단에 모아 보여 주기 때문에 도배 걱정이 없어요.
  \item 그룹 라이브: 채팅 오른쪽 위 "+" > [그룹 라이브]로 이 그룹만 볼 수 있는 라이브를 시작할 수 있습니다. 사내 교육, 가족 모임, 제품 내부 테스트에 적합합니다.
  \item 그룹 접기: 그룹 목록에서 그룹 이름을 왼쪽으로 스와이프 > [접기]하면 중요하지 않은 그룹을 "접은 그룹"으로 정리하고, 나를 @할 때만 알림이 옵니다.
  \item 그룹장/관리자만 @everyone: 멤버가 500명을 넘으면 일반 멤버는 전체 호출을 할 수 없어 방해를 줄일 수 있습니다.
\end{itemize}

\booksubsection{위챗 모멘트 권한 세밀 관리:}

\begin{itemize}
  \item 공개 범위 지정: 게시 시 "누가 볼 수 있나요"를 눌러 "일부 공개" 또는 "특정인 제외"를 선택하면, 미리 만든 태그 그룹(예: "동료", "가족", "고객") 기준으로 정확히 공유할 수 있습니다.
  \item 친구가 볼 수 있는 범위: [나] > [설정] > [친구 권한] > [모멘트] > [친구가 모멘트를 볼 수 있는 범위]에서 "최근 3일/1개월/6개월/전체" 중 선택할 수 있습니다. 프라이버시와 소셜의 균형을 위해 "6개월"을 권장합니다.
  \item 특정 모멘트 차단: 모멘트 게시물을 길게 눌러 [그의 모멘트 보지 않기] 또는 [팔로우 취소]를 선택하면 상대는 알 수 없습니다.
  \item 모멘트 광고 관리: [나] > [설정] > [WeChat 정보] > [개인정보 및 권한] > [맞춤 광고 관리]에서 행동 기반 맞춤 추천을 끌 수 있습니다.
  \item 낯선 사람의 모멘트 10개 열람: [친구 권한] > [모멘트] > [낯선 사람에게 모멘트 10개 공개]를 끄면 모르는 사람이 과거 게시물을 훑어보는 걸 막을 수 있습니다.
\end{itemize}

\booksection{영상 통화 비교}

영상 통화 도구마다 강점이 다릅니다. 상황에 맞는 도구를 고르면 소통이 훨씬 매끄럽습니다.

\begin{itemize}
  \item WeChat 영상: 전국민 커버리지, 남녀노소 사용. 가족·지인 일상 연락과 빠른 소통에 적합합니다. 최대 20명 영상(미니프로그램 "그룹 영상" 필요)을 지원하지만, 네트워크가 나쁘면 화질이 크게 떨어집니다.
  \item FaceTime: 애플 생태계 전용. 화질이 뛰어나고 종단 간 암호화를 지원합니다. iOS 15부터 "링크"로 안드로이드/Web 사용자를 초대할 수 있지만, 호스트는 애플 기기여야 합니다. 가족 간 기기 초월 소통에 좋습니다.
  \item Zoom: 국제 회의의 표준. 최대 100명 무료 40분 회의를 지원합니다. 화면 공유, 가상 배경, 분과 토론실 등 기능이 잘 갖춰져 있어 글로벌/원격 협업에 적합합니다.
  \item Tencent Meeting: 중국 오피스의 표준. WeChat과 깊게 연동되어 QR만 스캔하면 입장합니다. 300명 회의(유료), 클라우드 녹화, 실시간 자막, 회의록 기능이 강력해 사내 교육과 대형 회의에 좋습니다.
  \item 사용 장면 추천: 1:1 일상 연락은 WeChat, 애플 사용자끼리는 FaceTime, 공식 비즈니스 회의는 Zoom 또는 Tencent Meeting(중국 내부는 Tencent, 국경 간은 Zoom).
\end{itemize}

\booksubsection{네트워크 대역폭 및 데이터 사용 참고:}

\begin{itemize}
  \item WeChat 영상 통화: 약 500Kbps\textasciitilde{}1Mbps 필요. 표준 화질은 분당 약 3\textasciitilde{}5MB, 고화질은 약 8\textasciitilde{}12MB를 소비합니다.
  \item FaceTime: SD는 약 1Mbps, 720p는 약 2Mbps, 1080p는 4Mbps 이상. 분당 데이터는 약 5\textasciitilde{}15MB입니다.
  \item Zoom/Tencent Meeting: 720p는 약 1.5\textasciitilde{}3Mbps, 1080p는 3\textasciitilde{}5Mbps. 시간당 약 300\textasciitilde{}800MB를 쓰며, 다자 회의에서는 진행자의 데이터가 더 큽니다.
  \item 팁: 긴 영상 통화는 가능하면 WiFi를 사용하세요. 모바일 데이터에서는 앱 설정에서 "고화질 영상"을 꺼 데이터를 아낄 수 있습니다.
\end{itemize}

\booksubsection{영상 통화 화질 최적화 팁:}

\begin{itemize}
  \item 조명: 창문을 등지고 광원을 향하세요(자연광이 최고). 역광은 얼굴을 어둡게 만듭니다.
  \item 각도: 카메라를 눈보다 살짝 높이면 턱선이 정리됩니다. 가로 화면 통화는 양손으로 잡으면 더 안정적입니다.
  \item 오디오: 이어폰을 쓰면 에코와 하울링을 줄일 수 있습니다. 조용한 곳에서는 "스피커폰"을 켜지 마세요.
  \item 네트워크: WiFi 우선. 4G/5G 영상은 신호 약화를 피해 창가에 가까이 있는 게 좋습니다. 다자 회의는 다른 화면을 끄고 발화자만 남기세요.
  \item 배경 흐리게: Tencent Meeting, Zoom 모두 AI 배경 흐림을 지원해 어수선한 배경을 가릴 수 있습니다. 배경색과 비슷한 옷은 "크로마키 사고"를 부를 수 있으니 단색 옷을 입으세요.
\end{itemize}

\begin{tcbinfo}{꿀팁}
네트워크가 끊길 때는 바로 끊지 말고 영상만 끄고 음성으로 전환하는 것이 예의입니다. 중요한 회의 전에는 Tencent Meeting의 “네트워크 테스트”로 대역폭을 미리 확인하세요.

\end{tcbinfo}

\booksection{이메일: 휴대폰에서 설정하고 쓰기}

\begin{itemize}
  \item 이메일 계정 추가: iOS는 설정 > 메일 > 계정 > 계정 추가, 안드로이드는 보통 "이메일" 앱이 있습니다.
  \item 사용 팁: 이메일 계정마다 알림 방식을 다르게 설정하세요. 스와이프 제스처로 빠르게 보관·삭제·별표 표시하는 습관도 좋습니다.
\end{itemize}

\booksubsection{IMAP vs POP3 vs Exchange 프로토콜 차이:}

\begin{itemize}
  \item POP3: 서버에서 메일을 내려받으면 기본적으로 원본을 삭제하고 로컬에만 남깁니다. 기기 한 대만 쓰고 오프라인 열람이 필요할 때 적합합니다. 단, 기기를 바꾸면 과거 메일을 볼 수 없어 지금은 권장하지 않습니다.
  \item IMAP: 메일은 항상 서버에 있고 로컬은 캐시만 씁니다. 휴대폰/PC/태블릿 등 여러 기기의 상태가 일관됩니다. 한 곳에서 읽으면 모두 읽음 처리. 현재 주류로 권장되는 프로토콜입니다.
  \item Exchange(ActiveSync 포함): 마이크로소프트 프로토콜로, 메일뿐 아니라 캘린더·연락처·할 일까지 동기화합니다. 기업 메일의 1순위입니다. Outlook, NetEase 기업 메일, Tencent 기업 메일이 모두 지원합니다.
  \item 선택 가이드: 개인 메일은 IMAP, 회사 메일은 Exchange 우선, 그다음 IMAP. 특별한 아카이브 목적이 아니면 POP3는 피하세요.
\end{itemize}

\booksubsection{기업 메일 설정 방법:}

\begin{enumerate}
  \item iOS: [설정] > [메일] > [계정] > [계정 추가] > "Microsoft Exchange" 또는 "기타" 선택.
  \item 안드로이드: 기본 "이메일" 앱이나 Outlook 앱에서 "계정 추가"를 선택하고 전체 이메일 주소와 비밀번호를 입력합니다.
  \item 서버 입력(Exchange 기준): 서버 주소는 보통 "mail.회사도메인.com" 또는 "outlook.office365.com"이고, 도메인에는 회사 약칭을 넣습니다.
  \item "서버를 확인할 수 없습니다"가 뜨면 수동으로 입력하세요: 수신 서버(IMAP) 포트 993, 암호화 SSL. 발신 서버(SMTP) 포트 465 또는 587, 암호화 TLS.
  \item 로그인 후 [계정 설정]에서 "푸시"(Push)를 켜 실시간으로 메일을 받고, 영업 비밀 보호를 위해 "메일 미리보기"는 끄는 것을 권장합니다.
\end{enumerate}

\begin{tcbinfo}{알아두기}
기업 메일은 보통 관리자가 “클라이언트 승인 비밀번호”나 “IMAP/SMTP 승인 코드”를 열어 줘야 하며, 웹 로그인 비밀번호를 그대로 쓸 수 없습니다. 로그인에 실패하면 웹 메일 설정에서 승인 코드를 발급받아 휴대폰에 넣으세요.

\end{tcbinfo}

\begin{tcbpractice}{실습: WeChat 채팅 기록을 PC에 한 번 백업하기}
\begin{enumerate}
  \item PC에서 백업 시작: PC용 WeChat을 열고 왼쪽 아래 메뉴 > [이전 및 백업] > [백업 및 복원]으로 들어갑니다.
  \item 컴퓨터로 백업 선택: "채팅 기록을 컴퓨터에 백업"을 클릭합니다.
  \item 휴대폰에서 백업할 내용을 고르고 완료될 때까지 기다립니다.
\end{enumerate}

\end{tcbpractice}
